
\documentclass{article}
\usepackage{colortbl}
\usepackage{makecell}
\usepackage{multirow}
\usepackage{supertabular}

\begin{document}

\newcounter{utterance}

\centering \large Interaction Transcript for game `imagegame', experiment `compact\_grids', episode 7 with NEW\_base\_planning\_clp\_chat.
\vspace{24pt}

{ \footnotesize  \setcounter{utterance}{1}
\setlength{\tabcolsep}{0pt}
\begin{supertabular}{c@{$\;$}|p{.15\linewidth}@{}p{.15\linewidth}p{.15\linewidth}p{.15\linewidth}p{.15\linewidth}p{.15\linewidth}}
    \# & $\;$A & \multicolumn{4}{c}{Game Master} & $\;\:$B\\
    \hline

    \theutterance \stepcounter{utterance}  
    & & \multicolumn{4}{p{0.6\linewidth}}{
        \cellcolor[rgb]{0.9,0.9,0.9}{
            \makecell[{{p{\linewidth}}}]{
                \texttt{\tiny{[P1$\langle$GM]}}
                \texttt{Let us do the following task. The goal is to fill an empty grid that looks like this:} \\
\\ 
\texttt{▢ ▢ ▢ ▢ ▢} \\
\texttt{▢ ▢ ▢ ▢ ▢} \\
\texttt{▢ ▢ ▢ ▢ ▢} \\
\texttt{▢ ▢ ▢ ▢ ▢} \\
\texttt{▢ ▢ ▢ ▢ ▢} \\
\\ 
\texttt{A filled grid below is 5 by 5 and can look like this:} \\
\\ 
\texttt{▢ ▢ ▢ ▢ ▢} \\
\texttt{▢ ▢ E ▢ ▢} \\
\texttt{▢ ▢ ▢ ▢ ▢} \\
\texttt{▢ ▢ ▢ ▢ ▢} \\
\texttt{X X X X X} \\
\\ 
\texttt{I want you to describe this grid to me, step by step. You don't need to describe the empty squares, which are denoted with "▢". Only describe the location of letters in the grid. Then you wait for me to say "What is your next command?", and then you continue with the next step. Take the size of the grid into consideration while giving commands. You can give only one command at each turn but it can refer to single or multiple cells, rows, columns or even the whole grid at once. When you have described everything, you say "Command: DONE". Only give the command and omit everything else.} \\
\\ 
\texttt{For the filled grid above, here are the example steps.} \\
\\ 
\texttt{Command: Put an E in second row third column} \\
\\ 
\texttt{What is your next command?} \\
\texttt{Command: Fill the last row with X} \\
\\ 
\texttt{What is your next command?} \\
\texttt{Command: DONE} \\
\\ 
\texttt{Ok. Please do this for the following example, which is a 5 by 5 grid.} \\
\texttt{Z Z Z Z Z} \\
\texttt{▢ Z Z Z Z} \\
\texttt{▢ ▢ Z Z Z} \\
\texttt{▢ ▢ ▢ Z Z} \\
\texttt{▢ ▢ ▢ ▢ Z} \\
\texttt{What is your next command?} \\
            }
        }
    }
    & \\ \\

    \theutterance \stepcounter{utterance}  
    & \multicolumn{4}{p{0.45\linewidth}}{
        \cellcolor[rgb]{0.95,0.95,0.95}{
            \makecell[{{p{\linewidth}}}]{
                \texttt{\tiny{[P1$|$internal]}}
                \texttt{{[}} \\
\texttt{  \{} \\
\texttt{    "role": "user",} \\
\texttt{    "content": "Let us do the following task. The goal is to fill an empty grid that looks like this:\textbackslash{}n\textbackslash{}n\textbackslash{}u25a2 \textbackslash{}u25a2 \textbackslash{}u25a2 \textbackslash{}u25a2 \textbackslash{}u25a2\textbackslash{}n\textbackslash{}u25a2 \textbackslash{}u25a2 \textbackslash{}u25a2 \textbackslash{}u25a2 \textbackslash{}u25a2\textbackslash{}n\textbackslash{}u25a2 \textbackslash{}u25a2 \textbackslash{}u25a2 \textbackslash{}u25a2 \textbackslash{}u25a2\textbackslash{}n\textbackslash{}u25a2 \textbackslash{}u25a2 \textbackslash{}u25a2 \textbackslash{}u25a2 \textbackslash{}u25a2\textbackslash{}n\textbackslash{}u25a2 \textbackslash{}u25a2 \textbackslash{}u25a2 \textbackslash{}u25a2 \textbackslash{}u25a2\textbackslash{}n\textbackslash{}nA filled grid below is 5 by 5 and can look like this:\textbackslash{}n\textbackslash{}n\textbackslash{}u25a2 \textbackslash{}u25a2 \textbackslash{}u25a2 \textbackslash{}u25a2 \textbackslash{}u25a2\textbackslash{}n\textbackslash{}u25a2 \textbackslash{}u25a2 E \textbackslash{}u25a2 \textbackslash{}u25a2\textbackslash{}n\textbackslash{}u25a2 \textbackslash{}u25a2 \textbackslash{}u25a2 \textbackslash{}u25a2 \textbackslash{}u25a2\textbackslash{}n\textbackslash{}u25a2 \textbackslash{}u25a2 \textbackslash{}u25a2 \textbackslash{}u25a2 \textbackslash{}u25a2\textbackslash{}nX X X X X\textbackslash{}n\textbackslash{}nI want you to describe this grid to me, step by step. You don't need to describe the empty squares, which are denoted with \textbackslash{}"\textbackslash{}u25a2\textbackslash{}". Only describe the location of letters in the grid. Then you wait for me to say \textbackslash{}"What is your next command?\textbackslash{}", and then you continue with the next step. Take the size of the grid into consideration while giving commands. You can give only one command at each turn but it can refer to single or multiple cells, rows, columns or even the whole grid at once. When you have described everything, you say \textbackslash{}"Command: DONE\textbackslash{}". Only give the command and omit everything else.\textbackslash{}n\textbackslash{}nFor the filled grid above, here are the example steps.\textbackslash{}n\textbackslash{}nCommand: Put an E in second row third column\textbackslash{}n\textbackslash{}nWhat is your next command?\textbackslash{}nCommand: Fill the last row with X\textbackslash{}n\textbackslash{}nWhat is your next command?\textbackslash{}nCommand: DONE\textbackslash{}n\textbackslash{}nOk. Please do this for the following example, which is a 5 by 5 grid.\textbackslash{}nZ Z Z Z Z\textbackslash{}n\textbackslash{}u25a2 Z Z Z Z\textbackslash{}n\textbackslash{}u25a2 \textbackslash{}u25a2 Z Z Z\textbackslash{}n\textbackslash{}u25a2 \textbackslash{}u25a2 \textbackslash{}u25a2 Z Z\textbackslash{}n\textbackslash{}u25a2 \textbackslash{}u25a2 \textbackslash{}u25a2 \textbackslash{}u25a2 Z\textbackslash{}nWhat is your next command?"} \\
\texttt{  \},} \\
\texttt{  \{} \\
\texttt{    "role": "user",} \\
\texttt{    "content": "You are a skilled planner. Above you see the game history. What would be the best next action? Openly reason about the next step, then produce your next action exactly as the game expects it preceded by the separator 'ACT:' on its own line. If multiple actions are required, write ACT: first and then all actions."} \\
\texttt{  \}} \\
\texttt{{]}} \\
            }
        }
    }
    & & \\ \\

    \theutterance \stepcounter{utterance}  
    & \multicolumn{4}{p{0.45\linewidth}}{
        \cellcolor[rgb]{0.95,0.95,0.95}{
            \makecell[{{p{\linewidth}}}]{
                \texttt{\tiny{[P1$|$internal]}}
                \texttt{ACT: The grid has Zs along the main diagonal from the top{-}left to the bottom{-}right. The first row is completely filled with Zs, and each subsequent row has one fewer Z, starting from the left and moving right. The pattern is that Zs are placed in positions (1,1), (2,2), (3,3), (4,4), (5,5), and also in the entire first row. However, the first row is fully Z, which includes (1,1) to (1,5). The second row has Z in (2,2) to (2,5), but (2,1) is empty. The third row has Z in (3,3) to (3,5), and so on. The pattern is that for row i, columns from i to 5 are Z, and the rest are empty. So the filled Zs are: (1,1), (1,2), (1,3), (1,4), (1,5), (2,2), (2,3), (2,4), (2,5), (3,3), (3,4), (3,5), (4,4), (4,5), (5,5). But the first row is all Zs, which is a full row. The rest of the Zs are on the diagonal and to the right. So the pattern is: for} \\
            }
        }
    }
    & & \\ \\

    \theutterance \stepcounter{utterance}  
    & \multicolumn{4}{p{0.6\linewidth}}{
        \cellcolor[rgb]{0.9,0.9,0.9}{
            \makecell[{{p{\linewidth}}}]{
                \texttt{\tiny{[P1$\rangle$GM]}}
                \texttt{The grid has Zs along the main diagonal from the top{-}left to the bottom{-}right. The first row is completely filled with Zs, and each subsequent row has one fewer Z, starting from the left and moving right. The pattern is that Zs are placed in positions (1,1), (2,2), (3,3), (4,4), (5,5), and also in the entire first row. However, the first row is fully Z, which includes (1,1) to (1,5). The second row has Z in (2,2) to (2,5), but (2,1) is empty. The third row has Z in (3,3) to (3,5), and so on. The pattern is that for row i, columns from i to 5 are Z, and the rest are empty. So the filled Zs are: (1,1), (1,2), (1,3), (1,4), (1,5), (2,2), (2,3), (2,4), (2,5), (3,3), (3,4), (3,5), (4,4), (4,5), (5,5). But the first row is all Zs, which is a full row. The rest of the Zs are on the diagonal and to the right. So the pattern is: for} \\
            }
        }
    }
    & & \\ \\

    \theutterance \stepcounter{utterance}  
    & & & \multicolumn{2}{p{0.3\linewidth}}{
        \cellcolor[rgb]{0.9,0.9,0.9}{
            \makecell[{{p{\linewidth}}}]{
                \texttt{\tiny{[GM$|$GM]}}
                \texttt{Invalid instruction format} \\
            }
        }
    }
    & & \\ \\

\end{supertabular}
}

\end{document}
